\section{The Lindhard function at zero temperature}\label{app:D.8}

The density response of free electrons at $T = 0$ is a ratio of logarithms in $\omega$ and $q$, first written down by Lindhard~\cite{lindhard1954properties}. \Cref{sec:4.5} uses it for the valence electrons of silicon, where it fixes the static screening, the plasmon and the continuum of electron--hole pairs. The derivation follows Ref.~\cite[Chap.~4]{giuliani2005quantum} in the notation of \cref{sec:4.2,sec:4.7}.

\subsection*{The angular and radial integrals}

The independent-electron response is \cref{eq:4.7-bubble} continued to the real axis, $R^0(\vect{q},\omega) = \Vol^{-1}\sum_{\vect{k}}[n_{\vect{k}} - n_{\vect{k}+\vect{q}}]/(\omega + \ii0^+ + \xi_{\vect{k}} - \xi_{\vect{k}+\vect{q}})$, with the spin in the label. We keep the frequency as a complex number $z$ with $\Imag z > 0$, so that every expression below is analytic, and set $z = \omega + \ii0^+$ only at the end. At $T = 0$ the occupation is the Fermi sea of \cref{eq:2.5-fermi}, $n_{\vect{k}} = \Theta(k_{\mathrm{F}} - k)$. In the term with $n_{\vect{k}+\vect{q}}$, relabel $\vect{k} \to -\vect{k} - \vect{q}$. The occupation becomes $n_{\vect{k}}$ and the denominator becomes $z + \xi_{\vect{k}+\vect{q}} - \xi_{\vect{k}}$, since $\xi$ depends on $|\vect{k}|$ only. With the two spins and $\sum_{\vect{k}} \to \Vol\int\dd^3k/(2\pi)^3$,
\begin{equation*}
  R^0(q,z) = 2\int_{k < k_{\mathrm{F}}}\frac{\dd^3k}{(2\pi)^3}\left[\frac{1}{z - \Delta_{\vect{k}}} - \frac{1}{z + \Delta_{\vect{k}}}\right], \qquad
  \Delta_{\vect{k}} = \frac{(\vect{k} + \vect{q})^2 - k^2}{2m_e} = \frac{kqu}{m_e} + \frac{q^2}{2m_e} ,
\end{equation*}
with $u$ the cosine of the angle between $\vect{k}$ and $\vect{q}$. The resonant term, with $z - \Delta_{\vect{k}}$, is an electron leaving the sea at $\vect{k}$ for $\vect{k} + \vect{q}$, the transition that absorbs energy; the term with $z + \Delta_{\vect{k}}$ is its antiresonant partner of \cref{eq:4.1-lehmann}. Only states inside the sphere appear, and the Pauli principle has moved into a cancellation between the two terms. A transition from $\vect{k}$ to an occupied $\vect{k} + \vect{q}$ appears in the first term at $\vect{k}$ and, with the opposite sign, in the second term at $-\vect{k} - \vect{q}$, which is also inside the sphere, and the two cancel.

The variables that make the result simple measure frequency and momentum in units set by the Fermi surface,
\begin{equation*}
  \nu = \frac{z}{qv_{\mathrm{F}}}, \qquad y = \frac{q}{2k_{\mathrm{F}}}, \qquad \tilde{k} = \frac{k}{k_{\mathrm{F}}}, \qquad \nu_\pm = \nu \pm y ,
\end{equation*}
with $v_{\mathrm{F}} = k_{\mathrm{F}}/m_e$. Then $z \mp \Delta_{\vect{k}} = qv_{\mathrm{F}}(\nu \mp y \mp \tilde{k} u)$. The angular integrals are elementary, $\int_{-1}^1\dd u\,(w - \tilde{k} u)^{-1} = \tilde{k}^{-1}\ln[(w + \tilde{k})/(w - \tilde{k})]$ for complex $w$ off the real segment $[-\tilde{k}, \tilde{k}]$, with the principal branch of the logarithm. For $\Imag w > 0$ the ratio $(w + \tilde{k})/(w - \tilde{k}) = 1 + 2\tilde{k}/(w - \tilde{k})$ lies in the lower half plane, away from the cut of the principal branch, so the logarithm is analytic there. With $\dd^3k = 2\pi k^2\,\dd k\,\dd u$,
\begin{equation*}
  R^0(q,z) = \frac{m_ek_{\mathrm{F}}^2}{2\pi^2q}\int_0^1\dd\tilde{k}\;\tilde{k}\left[\ln\frac{\nu_- + \tilde{k}}{\nu_- - \tilde{k}} - \ln\frac{\nu_+ + \tilde{k}}{\nu_+ - \tilde{k}}\right] .
\end{equation*}
The radial integral is the one of the exchange energy in \cref{app:D.4}. Its antiderivative, $\tfrac12(\tilde{k}^2 - w^2)\ln[(w + \tilde{k})/(w - \tilde{k})] + w\tilde{k}$, holds for complex $w$, since the derivative is taken without using reality. Evaluated between $0$ and $1$ it gives
\begin{equation*}
  \mathcal{I}(w) = \int_0^1\dd\tilde{k}\;\tilde{k}\ln\frac{w + \tilde{k}}{w - \tilde{k}} = w + \frac{1 - w^2}{2}\,\ell(w), \qquad \ell(w) = \ln\frac{w + 1}{w - 1} ,
\end{equation*}
and $R^0 = (m_ek_{\mathrm{F}}^2/2\pi^2q)[\mathcal{I}(\nu_-) - \mathcal{I}(\nu_+)]$. Since $\nu_- - \nu_+ = -2y = -q/k_{\mathrm{F}}$, the linear terms of $\mathcal{I}$ give a constant. Collecting, with the density of states $\mathcal{N}_{\mathrm{F}} = m_ek_{\mathrm{F}}/2\pi^2$ of one spin at the Fermi level (\cref{eq:2.5-fermi}), we arrive at the Lindhard function,
\begin{equation}
  R^0(q,z) = -2\mathcal{N}_{\mathrm{F}}\left\{\frac12 + \frac{k_{\mathrm{F}}}{4q}\Bigl[(1 - \nu_+^2)\,\ell(\nu_+) - (1 - \nu_-^2)\,\ell(\nu_-)\Bigr]\right\}, \qquad
  \nu_\pm = \frac{z}{qv_{\mathrm{F}}} \pm \frac{q}{2k_{\mathrm{F}}} .
  \label{eq:D.8-Rz}
\end{equation}
It is analytic in the upper half plane of $z$, as a retarded function must be (\cref{app:C.5}), and it depends on $q$ and $z$ only through $\nu$ and $y$, apart from the prefactor $2\mathcal{N}_{\mathrm{F}}$.

\subsection*{Real and imaginary parts}

On the real axis, $z = \omega + \ii0^+$ with $\omega > 0$, and both $\nu_\pm$ acquire a positive imaginary part. By \cref{eq:C.8-branch}, $\ell(\nu_\pm + \ii0^+) = \ln|(\nu_\pm + 1)/(\nu_\pm - 1)| - \ii\pi\,\Theta(1 - |\nu_\pm|)$. The real part of \cref{eq:D.8-Rz} keeps the form of the equation with every logarithm replaced by the logarithm of the modulus, now with $\nu_\pm = \omega/qv_{\mathrm{F}} \pm q/2k_{\mathrm{F}}$ real,
\begin{equation}
  \Real R^0(q,\omega) = -2\mathcal{N}_{\mathrm{F}}\left\{\frac12 + \frac{k_{\mathrm{F}}}{4q}\left[(1 - \nu_+^2)\ln\Bigl|\frac{\nu_+ + 1}{\nu_+ - 1}\Bigr| - (1 - \nu_-^2)\ln\Bigl|\frac{\nu_- + 1}{\nu_- - 1}\Bigr|\right]\right\} ,
  \label{eq:D.8-Re}
\end{equation}
and the imaginary part collects the steps,
\begin{equation}
  \Imag R^0(q,\omega) = -\frac{\pi k_{\mathrm{F}}\mathcal{N}_{\mathrm{F}}}{2q}\Bigl[(1 - \nu_-^2)\,\Theta(1 - |\nu_-|) - (1 - \nu_+^2)\,\Theta(1 - |\nu_+|)\Bigr] .
  \label{eq:D.8-Im}
\end{equation}
For $\omega < 0$ the imaginary part changes sign, since $R^0(q,-\omega) = R^0(q,\omega)^*$ by \cref{eq:4.1-reality} for a gas with $R^0(-\vect{q}) = R^0(\vect{q})$.

The steps divide the $(q, \omega)$ plane into a region where both are on, a region where only the step of $\nu_-$ is on, and a region outside the continuum. Since $\omega > 0$, $\nu_+ > \nu_-$ and $\nu_+ > y > 0$. The condition $|\nu_+| < 1$ is $\omega < qv_{\mathrm{F}} - q^2/2m_e$, which needs $q < 2k_{\mathrm{F}}$. The condition $|\nu_-| < 1$ is $|m_e\omega/q - q/2| < k_{\mathrm{F}}$, which is the band $\omega_- < \omega < \omega_+$ with $\omega_+ = qv_{\mathrm{F}} + q^2/2m_e$ and $\omega_-$ equal to $q^2/2m_e - qv_{\mathrm{F}}$ for $q > 2k_{\mathrm{F}}$ and to zero otherwise. Below the line $\omega = qv_{\mathrm{F}} - q^2/2m_e$ both steps are on. The bracket is then $\nu_+^2 - \nu_-^2 = 4\nu y = 2m_e\omega/k_{\mathrm{F}}^2$, and
\begin{equation*}
  \Imag R^0 = -\frac{m_e^2\,\omega}{2\pi q} \qquad (\omega < qv_{\mathrm{F}} - q^2/2m_e) .
\end{equation*}
Above it, up to $\omega_+$, only the step of $\nu_-$ is on, and with $k_{\mathrm{F}}^2(1 - \nu_-^2) = k_{\mathrm{F}}^2 - (m_e\omega/q - q/2)^2$,
\begin{equation*}
  \Imag R^0 = -\frac{m_e}{4\pi q}\Bigl[k_{\mathrm{F}}^2 - \Bigl(\frac{m_e\omega}{q} - \frac{q}{2}\Bigr)^2\Bigr] \qquad (\text{only } |\nu_-| < 1) .
\end{equation*}
Outside the band $\omega_- < \omega < \omega_+$ the imaginary part vanishes: no electron--hole pair has that energy at that momentum. The two forms are $-\tfrac12$ times the linear and parabolic forms of the structure factor $S^0$ in \cref{we:4-gas}, which is the fluctuation--dissipation theorem \cref{eq:4.2-fdt} at $T = 0$, $S^0 = -2\Imag R^0$ for $\omega > 0$. The two calculations are independent: \cref{we:4-gas} counted the transitions with a delta function, and here they appear as the jump of a logarithm across its cut.

\subsection*{Static and small-$q$ limits}

\emph{Static response.} At $\omega = 0$, $\nu_\pm = \pm y$. The logarithm of the modulus is odd, $\ln|(1 - y)/(-1 - y)| = -\ln|(1 + y)/(1 - y)|$, so the two terms in the bracket of \cref{eq:D.8-Re} add, and $(k_{\mathrm{F}}/4q)\cdot2 = 1/4y$:
\begin{equation}
  R^0(q,0) = -2\mathcal{N}_{\mathrm{F}}\left[\frac12 + \frac{1 - y^2}{4y}\ln\Bigl|\frac{1 + y}{1 - y}\Bigr|\right] = -2\mathcal{N}_{\mathrm{F}}\,F_{\mathrm{L}}\!\left(\frac{q}{2k_{\mathrm{F}}}\right) ,
  \label{eq:D.8-static}
\end{equation}
with $F_{\mathrm{L}}$ of \cref{eq:C.8-FL}.

\emph{Small momentum at fixed $\omega/qv_{\mathrm{F}}$.} For $y \to 0$ at fixed $\nu$, $\mathcal{I}(\nu - y) - \mathcal{I}(\nu + y) = -2y\,\mathcal{I}'(\nu) + O(y^3)$. The derivative is $\mathcal{I}'(w) = 1 - w\,\ell(w) + \tfrac12(1 - w^2)\ell'(w)$, and with $\ell'(w) = -2/(w^2 - 1)$ the last term is $1$, so $\mathcal{I}'(w) = 2 - w\,\ell(w)$. Hence
\begin{equation}
  R^0(q,z) = -2\mathcal{N}_{\mathrm{F}}\left[1 - \frac{\nu}{2}\ln\frac{\nu + 1}{\nu - 1}\right] + O(y^2), \qquad \nu = \frac{z}{qv_{\mathrm{F}}} .
  \label{eq:D.8-smallq}
\end{equation}
This is the response of a Fermi surface to a long-wavelength disturbance, the semiclassical limit in which the recoil $q^2/2m_e$ is negligible against $qv_{\mathrm{F}}$. For $0 < \omega < qv_{\mathrm{F}}$ the imaginary part is $-\pi\mathcal{N}_{\mathrm{F}}\,\omega/(qv_{\mathrm{F}})$, which equals $-m_e^2\omega/2\pi q$, the linear form above; above $qv_{\mathrm{F}}$ it vanishes.
